\documentclass[10pt]{article}
\usepackage[margin=0.75in]{geometry}
\usepackage{amsmath,amssymb,amsthm,microtype}
\newtheorem{theorem}{Theorem}
\newcommand{\R}{\mathcal R}
\newcommand{\A}{\mathcal A}
\begin{document}
\begin{center}
{\Large An Inverse Excursion That Lowers Cycle Charge}\par\medskip
Collatz proof project\quad---\quad September 5, 2026
\end{center}
Let $T$ be the shortcut Collatz map, $j_t(n)$ its odd-step count,
$\R(n)$ eventual reachability of one, and
$\A(u)=[\R(3u+2)\Rightarrow\R(27u+20)]$.
The previous growth bound excludes lowering cycle charge using only forward
steps and auxiliary bridges $x\mapsto9x+2$. We now give a concrete signed
excursion outside that bound. Its recursive transfer premise remains explicit.

\begin{theorem}[Kernel-checked inverse bridge]
For every natural $Q$,
\begin{align*}
T^{65}(3(15+2^{65}Q)+2)&=2+3^{39}Q,\\
T^{38}(27(15+2^{65}Q)+20)&=2+3^{22}2^{27}Q.
\end{align*}
Consequently, assuming $\A(15+2^{65}Q)$,
\[
 \R(2+3^{39}Q)\ \Longrightarrow\ \R(2+3^{22}2^{27}Q).
\]
\end{theorem}
\begin{proof}
The base computations are
$T^{65}(47)=2$, $j_{65}(47)=38$,
$T^{38}(425)=2$, and $j_{38}(425)=19$.
Binary-cylinder transport gives the displayed identities.
Transport convergence backward through the first identity, apply the explicit
transfer premise, and transport forward through the second.
Lean checks both base paths, both identities for every $Q$, and the implication.
\end{proof}

\paragraph{Effect on the symbolic state.}
At a cycle state $(2,p,s)$, representing $2+3^p2^{D-s}Q$, reverse the
65-step path from 47, use its bridge $47\mapsto425$, and advance 38 steps.
The resulting state is $(2,p-17,s-27)$, provided the reversed lift is integral.
Thus the cycle charge $L(x,p,s)=2p-s+\mathbf1_{x=1}$ changes by
$2(-17)-(-27)=-7$. This does not contradict the forward-only theorem.

\paragraph{Admissibility in parameter induction.}
The original parameter is $U=u+2^D Q$.
For this excursion the recursive parameter is
\[
 v=15+3^{p-39}2^{D-s+65}Q.
\]
Sufficient conditions are $u>15$, $p\geq39$, $s\geq65$, and
$3^{p-39}\leq2^{s-65}$. They ensure integrality and $v<U$ for every $Q$.
Waiting two ordinary steps at base state 2 raises $p$ by one and $s$ by two;
the slope ratio is multiplied by $3/4$. Sufficiently long waiting therefore
satisfies the slope condition and the required minimum clock and exponent.
Waiting is also chosen so the target remains past its cycle entry when the
source clock is moved backward.

\paragraph{Conditional local completeness: a written construction.}
Fix $u>15$ and suppose both finite base partners reach the trivial cycle.
Their charge difference is an integer. Repeated $-7$ excursions can lower the
source charge below or to the target charge. Repeated $+1$ excursions
$2\mapsto20\to\cdots\to2$ then make the charges equal.
The latter uses seven ordinary steps with two odd steps and recursive base
parameter zero. Each finite excursion can be made admissible by waiting.
At a common clock, equal cycle charge determines both cycle phase and slope.
Choose $D$ at least the maximum clock visited, not merely the net final clock,
and advance both matched endpoints to $D$ if necessary.
This gives a conditional all-quotient transfer certificate on $u+2^D Q$.
The complete scheduling argument is not yet formalized in Lean; the displayed
inverse macro is formalized.

\paragraph{Finite verification.}
The implementation constructs certificates for all 985 parameters from 16
through 1000; the largest final depth is 412. Every generated certificate is
replayed at $Q=0,1,7$, checking reverse identities, all recursive decreases,
and equality of the final affine endpoints. Three tests pass. The saved report
contains the census and four representative paths, not a kernel-certified census.
The maximum-clock condition is tested explicitly on parameter 20, which needs
two negative excursions.

\paragraph{The unresolved global step.}
The construction uses convergence of the base target to determine its charge.
Assuming that for every parameter would already assume the convergence problem.
Thus local completeness for convergent base pairs is not universal induction
coverage. The new transition removes the earlier directional obstruction but
supplies no proof of Collatz. The three public declarations in
\texttt{Collatz.InverseCycleBridge} are included in the axiom audit.
No mathematical priority claim is made.
\end{document}
